% ch3 A 元激发的概念：多体理论概述 (原书页 96–103 / PDF p111–p118)

\chapter{元激发}

\section{元激发的概念：多体理论概述}

我在这门课程第一学期中所讲的一切，都建立在一个近似之上——而事实上，我们很容易使自己相信，这个近似在多数情况下是非常粗糙的。那个近似就是：固体中的电子彼此独立地运动，也独立于重离子实（ion cores）的运动——后者我们甚至未曾讨论过。充其量，我们只不过以平均的方式计入了电子之间的相互作用——也就是说，只计到 Hartree-Fock 理论的程度——除了简略地提到 Wigner-Pines 对关联能（correlation energy）的计算之外。

事实上，固体中的电子和离子确实相互作用得非常强烈，而且可以预期是一种急速涨落式的相互作用。有两个简单的例子值得注意。我们已经看到，关联能修正与碱金属的结合能（binding energy）相当，因此很难把它看作一个小效应。我们还注意到，在相当大的距离上，两个电荷之间的相互作用能必须用介电常数（dielectric constant）$\kappa$ 来折减：$\kappa:\ e^{2}/r \to e^{2}/\kappa r$。即便对绝缘体，典型固体的 $\kappa$ 值也在 3 到 20 的量级，因此相互作用由此发生的改变相当可观。

这两者都是真正的多体效应，完全处在 Hartree-Fock 理论的领域之外。在这份讲义的其余部分，我要做的正是讨论这些多体效应。遗憾的是，在这一领域中，除了最高深的专著之外似乎没有一本好的书，尽管其中许多概念并不困难；最好的或许要数 Pines 的一套讲义和重印文集 (58)。其他参考文献有 Landau-Lifshitz（文献 6 的第六章），以及更高深但在高深著作中最易读的 Thouless《Quantum Mechanics of Many-Body Systems》(59)；还有 Les Houches 讲义 (80)。

那么，使用简单的单电子概念，为什么竟能如此经常地给我们带来定性上、有时甚至是半定量上正确的结论呢？这个问题实际上有不少答案，而在接下来的几讲中，我们将主要关心的只是其中之一。不过，在这里把其中几个列举出来或许是有益的：（1）变分定理（variational theorem），（2）不相容原理（exclusion principle），（3）屏蔽（screening），（4）元激发（elementary excitations）的概念。

\subsection{变分定理}

第一个是变分定理，它几乎不言自明——但它只适用于基态能量。这就是说，虽然 Hartree-Fock 波函数对一个包含许多电子和离子的系统来说，并不是其真实波函数的一个十分准确的近似（对大系统而言，两者之间几乎没有任何交叠），但它的能量作为一个近似却远没有那么糟糕。

\subsection{不相容原理}

在电子系统的 Hartree-Fock 理论中，不相容原理在削减多体效应方面起着非常重要的作用。这一效应虽然许多年前就已在自由电子气体中得到认识，但把它明确地指出并加以形式化的，主要是 Weisskopf (61)，其后是处理核子情形的 Brueckner；而有限电子系统中的类似效应，则由 MacWeeny (62) 作了特别的强调。

在 2-A1 节中，我们曾把一个多电子系统的 Hartree-Fock 波函数写成如下形式：
\begin{equation*}
\Psi_{\mathrm{HF}}=\prod_{\substack{\text{occ. states}\\ n,\ \sigma}}^{N}c^{\dagger}_{n\sigma}\ \Psi_{\mathrm{vac}}
\end{equation*}
其中 $c^{\dagger}_{n\sigma}$ 在一组正交轨道 $\phi_{n\sigma}$ 上产生电子；这组轨道是该问题最好的单电子轨道，其含义是：我们消除了所有那种只使单个电子独自从一个占据态 $\phi_{n\sigma}$ 激发到一个未占据态 $\phi_{M\sigma}$（$M>N$）去的矩阵元。正如我们解释过的，这就是方程
$[\mathcal{H},\ c^{\dagger}_{n\sigma}]=E_{n}\,c^{\dagger}_{n\sigma}+\text{truly off-diagonal terms}$
的含义。

于是，能够发生的任何微扰都必须同时把两个电子从占据态 $n$、$n'$ 送到未占据态 $m$、$m'$ 中去；相应的矩阵元是
\begin{align*}
&\left(\,mm'\,\middle|\,e^{2}/r_{12}\,\middle|\,nn'\,\right)\\
&\qquad=\int dr_{1}\int dr_{2}\ \frac{2}{|r_{1}-r_{2}|}\ \phi^{*}_{m}(r_{1})\ \phi^{*}_{m'}(r_{2})\ \phi_{n}(r_{1})\ \phi_{n'}(r_{2})
\end{align*}

例如，微扰论对能量的最低阶修正是
\begin{equation*}
-\sum_{\substack{n,\ n'<N\\ m,\ m>N}}\left|\,\left(mm'\,\middle|\,e^{2}/r_{12}\,\middle|\,nn'\right)\,\right|^{2}\Big/\left(E_{m}+E_{m'}-E_{n}-E_{n'}\right)
\end{equation*}
（译注：求和限第二行原书印作“m, m>N”，第二个 m 漏印撇号，应为 $m, m'>N$。）

由于不相容原理，这些矩阵元很可能相当小。原因在于，态 $m$ 和 $m'$ 必须从与态 $n$ 和 $n'$ \emph{正交}的一个态子空间中选取，因而很可能与它们大不相同——要么在真实空间中占据不同的位置，要么相对它们急速变化，也就是说，处在动量空间的不同部分。于是这个积分通常出人意料地小。有一种很好的表述方式，就是指出：这个积分乃是电荷分布
\begin{equation*}
\rho_{mn}=\phi^{\dagger}_{m}(r_{1})\ \phi_{n}(r_{1})\qquad\text{with}\qquad\rho_{m'n'}=\phi^{\dagger}_{m'}(r)\ \phi_{n'}(r)
\end{equation*}
的 Coulomb 相互作用能。但由正交性，
\begin{equation*}
\int dr\ \rho_{mn}(r)=\int dr\ \rho_{m'n'}(r)=0
\end{equation*}
也就是说，这些电荷分布的符号正负交替，而且通常变化得相当快，所以它们之间的相互作用不很大。这就是 MacWeeny 的论证。

核物理学家则把它表述成：由于不相容原理，电子 $n$ 和 $n'$ 彼此散射所能进入的相空间大受限制。例如，在自由电子气体中，它们必须把彼此完全散射到费米球（Fermi sphere）之外，这就削弱了散射（图 29）。我们在研究费米液体（Fermi liquids）和磁性状态时，都还要回到这个话题。

注意，每一对可能的 $n$、$n'$ 所发生的散射都会是有限的，因此真实的 $\Psi_{\mathrm{g}}$ 相对于 $\Psi_{\mathrm{HF}}$ 来说，对每一个 $n$ 都有一个有限的修改。于是，当 $N\to\infty$ 时，由于 $\Psi_{\mathrm{HF}}$ 是 $N$ 个因子的乘积，而其中每个因子都被修改了一个有限的量，它与真实的 $\Psi_{\mathrm{g}}$ 的交叠 $\to 0$。

%FIG{29}: 图 29

\subsection{屏蔽}

当我们在电子气体中引入一个电荷——例如，我们可能正在研究其运动的那一个电子的电荷——时，其他电子当然会受它吸引而向它移动（若它为正），或者被它排斥而远离（若它为负）。例如在后一种情形中，这种运动会在原有 Hartree-Fock 密度的基础上留下一处电子电荷的亏缺，它等价于一定量的正电荷。于是，这个外来电子在任何距离上与其他电子相互作用时，就好像它带着一个较小的电荷——它已经被其他电子\emph{屏蔽}了。这一效应非常强地起着削减关联效应的作用。

\subsection{元激发的概念}

最后一个观念，也是我将要用若干页的篇幅来谈的观念，就是元激发的观念。这个概念同样是一个逐渐成长起来的观念，而且成长得如此缓慢，以至于人们几乎说不出它的创始人是谁；不过，对于“存在着这样一种普遍的处理方法”这一自觉的认识，或许 Landau［见 1941 年的文献 (63) 及其著作 (6)］的功劳比任何人都大。“元激发”这个名称就是他引入的。Fröhlich (64) 和 Pines (65) 使人们认识到开发利用量子场论方法的价值，这一做法把多体理论这整个课题开拓成了一门数学学科；而 Kohn (66) 和 Luttinger (67)——连同其他一些人，但他们的贡献超过任何他人——则确立了这样一个观念：元激发理论中的某些思想是可以给出严格证明的。但是，人们不应漏掉 Debye 的名字，他是元激发理论（Debye 声子理论）的第一个使用者；也不应漏掉 Wigner 的名字，他是第一个就关联电子运动这一主题进行推理的人。

元激发概念背后有两个思想。第一个思想是：基态的总结合能并不是一个十分重要的物理量，它与一个物理系统的行为没有多大关系。物理上\emph{重要的}，是较低激发态相对于基态的行为：也就是说，那些在相对较低的温度下或弱的外场中就容易受到激发的态。我们立刻会想到金属或半导体——其中的一切行为都由那些我们说成具有少数几个运动载流子的低激发态所决定；或者想到固体的弹性或热学性质——它们由少数几个我们称之为声子（phonons）的晶格波的存在所决定。因此，我们的兴趣往往集中于一个系统的整组低激发态上，把它视为该系统在物理上最基本的性质。

第二个思想是：低激发态往往——事实上几乎总是——具有特别简单的性质，与其他的态相比，能够以大得多的数学严格性和物理通透性来处理。让我解释一下其中的原因。

正如我一直强调的，尽管有我们一直在谈论的这些削减效应，真实基态 $\Psi_{\mathrm{g}}$ 的波函数和能量所受到的多体修正仍然非常大：特别是 $\Psi_{\mathrm{g}}$ 与单粒子近似 $\Psi_{\mathrm{HF}}$ 之间几乎没有什么关系，甚至毫无关系，因为 $\Psi_{\mathrm{g}}$ 在全部 $N$ 个粒子的坐标上都与后者不同。

不过，关于基态的群性质，我们通常倒是知道一些的。例如在绝缘体中，电子波函数具有平移群的完全对称性：$T\Psi_{\mathrm{g}}=\Psi_{\mathrm{g}}$。

为简单起见，让我们考虑系统恰好只多加了一个粒子的那些激发态。在 Hartree-Fock 近似下，给定晶体动量 $k$ 的最低这样的态，应当是那个多余电子处在最低空带中动量为 $k$ 的状态：
\begin{equation*}
\Psi_{k}=\left(c^{\,n'}_{k}\right)^{\dagger}\ \Psi_{\mathrm{HF}}
\end{equation*}
这当然完全不像一个本征态。然而，我们可以把动量为 $k$、多一个电子的最低本征态 $\Psi_{k}$，与\emph{真实的}基态本征态 $\Psi_{\mathrm{g}}$ 加以比较。会存在某个算符 $q^{\dagger}_{k}$，它代表这两个态之间的关系：
\begin{equation*}
\Psi_{k}=q^{\dagger}_{k}\ \Psi_{\mathrm{g}}
\end{equation*}
关于算符 $q^{\dagger}_{k}$，最重要的一点是：它必然只代表整个系统的一个非常小的位移。在我们正在讨论的简单情形中，我们将看到，事实上在 $N\sim 10^{23}$ 个电子中，只有动量为 $k$ 的那些电子才受到有限大小的扰动。在更复杂的情形中，也许每个电子都在某种意义上被移动了——比如在等离体振荡（plasma oscillation）中那样——但移动的量只是无穷小。

因此，就几乎所有电子而论，$\Psi_{k}$ 与 $\Psi_{\mathrm{g}}$ 几乎是相同的。例如，如果我们考虑一个由靠近中心值 $k_{\mathrm{o}}$ 的那些 $k$ 组成的波包（wave packet），
\begin{equation*}
\Psi_{\text{packet}}(k_{\mathrm{o}},\,r_{\mathrm{o}})=\int\exp\left[-ik\cdot r_{\mathrm{o}}-\frac{1}{2}\left(\frac{k-k_{\mathrm{o}}}{\Delta k}\right)^{2}\right]\Psi_{k}\ dk
\end{equation*}
我们就可以预期，波函数中由此产生的扰动将局限于 $r_{\mathrm{o}}$ 附近尺度为 $1/\Delta k$ 的区域内。如果波函数果真是 Hartree-Fock 的，这当然必然如此：$\Psi_{\text{packet}}$ 将会是
\begin{align*}
\left(\Psi_{\text{packet}}\right)_{\mathrm{HF}}&=\int e^{-ik\cdot r_{\mathrm{o}}}\ \exp\left[-\frac{1}{2}\left(\frac{k-k_{\mathrm{o}}}{\Delta k}\right)^{2}\right]c^{\dagger}_{k}\ \Psi_{\text{H-F}}\ dk\\
&=\int dr\left(\int e^{-ik\cdot(r_{\mathrm{o}}-r)}\ \exp\left[-\frac{1}{2}\left(\frac{k-k_{\mathrm{o}}}{\Delta k}\right)^{2}\right]dk\right)\Psi^{\dagger}(r)\ \Psi_{\mathrm{HF}}\\
&=\left(\int dr\ f_{\text{packet}}(r)\ \Psi^{\dagger}(r)\right)\Psi_{\mathrm{HF}}
\end{align*}
其中
\begin{equation*}
f_{\text{packet}}(r)\propto e^{\,ik_{\mathrm{o}}\cdot(r-r_{\mathrm{o}})}\exp\left[-\left(\frac{\Delta k}{2}\right)^{2}(r-r_{\mathrm{o}})^{2}\right]
\end{equation*}

在更一般的情形中，我们感到，波包的形成将造成对真实基态 $\Psi_{\mathrm{g}}$ 的一个局域性的扰动。于是，算符
\begin{equation*}
q^{\dagger}_{\text{packet}}(r_{\mathrm{o}},\,k_{\mathrm{o}})=\int dk\ \exp\left[-ik\cdot r_{\mathrm{o}}-\frac{1}{2}\left(\frac{k-k_{\mathrm{o}}}{\Delta k}\right)^{2}\right]q^{\dagger}_{k}
\end{equation*}
就是一个只产生局域扰动的算符。

于是，如果我们构造一个带有两个这种波包的波函数
\begin{align*}
&\Psi(k_{\mathrm{o}},\,r_{\mathrm{o}};\ k_{\mathrm{o}}',\,r_{\mathrm{o}}')\\
&\qquad=q^{\dagger}(k_{\mathrm{o}},r_{\mathrm{o}})\,q^{\dagger}(k_{\mathrm{o}}',r_{\mathrm{o}}')\ \Psi_{\mathrm{g}}\qquad\quad k_{\mathrm{o}}\neq k_{\mathrm{o}}',\ \ r_{\mathrm{o}}\neq r_{\mathrm{o}}'
\end{align*}
我们就可以预期，对 $r_{\mathrm{o}}$ 和 $r_{\mathrm{o}}'$ 的大多数可能取值，这两个波包彼此之间不会有多大干涉。于是，如果激发态 $\Psi_{k}$ 的激发能为 $E_{k}=(\Psi_{k}\,|\,\mathcal{H}\,|\,\Psi_{k})-E_{\mathrm{g}}$，我们就可以预期：含有两个这种激发的激发态，其能量是两个能量之和，准确到 $1/N$ 的量级（这正是两个波函数实际交叠的量级）。
\begin{equation*}
E(k_{\mathrm{o}},k_{\mathrm{o}}')=E_{k_{\mathrm{o}}}+E_{k_{\mathrm{o}}'}+\mathrm{O}\left(\frac{1}{N}\right)
\end{equation*}

这就是元激发概念背后的基本思想：在一个非常大的系统中，两个这样的激发将不会相互干涉，无论它们是像我在这里讨论的这种简单准粒子（quasi-particles），还是声子、激子（excitons），抑或是某种更加复杂的东西。因此，当存在的激发数目足够小时，系统的各种性质——特别是能量——将是无相互作用的元激发诸性质的线性叠加。这种情形下的另一个推论是：两个彼此不相互作用的 $q_{k}$ 将具有真实的费米子反对易规则。这一点证明起来并不容易，但它是真的：一般而言，我们总可以把这些 $q$ 取成具有费米子的对易规则。

这样，我们可以对元激发下一个初步的定义：\emph{动量为 $k$ 的元激发，就是那个从基态产生特定类型的、动量为 $k$ 的最低激发态的算符。}根据上面的推理，我们预期这些激发只在弱的程度上相互作用——事实上是 $1/N$ 的量级——因此系统中可以容纳相对大量的这种激发，从而按激发的绝对数目 $n$ 计，系统可以处在激发程度非常高的态，而我们仍然能够把这些激发当作一种近似无相互作用的独立粒子气体来处理。必须认识到，上述波包形成的想法并不限于我所讨论的那种特殊类型的激发态，它也可以应用于那些被适当选取、用以代表单个声子、自旋波（spin waves）等等的最低本征态。换一种说法：元激发的概念，是使系统的方程围绕真实基态——而不是围绕某个独立粒子近似——\emph{线性化}的一种方法。

在进一步讨论上述定义的种种限定与推广之前，最好先给出一些具体的例子。此外，我还想为我早先说过的那个论断提供支持：我们对固体和某些液体的物理分类，在很大程度上取决于它们所表现的元激发的类型。因此，我将用列表的形式写出元激发以及它们在其中显得重要的材料类型：元激发列在左边，它所表征的材料类型列在右边。

\emph{1. 自由电子载流子（free electronic charge carriers）}

\noindent\begin{tabular}{@{}p{0.58\textwidth}@{\hspace{2em}}p{0.36\textwidth}@{}}
a.~有能隙，能量极小值位于动量空间中一些孤立的点附近 & 绝缘体和半导体\\
b.~有能隙，能量极小值位于动量空间中的一个面上；电荷态反常 & 超导体\\
c.~无能隙；$E_{k}$ 在 $k$ 空间中的一个面上变为零 & 金属（正常的）\\
\end{tabular}

\emph{注意：}极化子（polarons）只不过是上列各项的进一步引申，唯有自陷极化子（self-trapped polaron）例外——在某些既非金属又非绝缘体的物质中，它可能成为占主导地位的激发。

\emph{2. 密度波（density waves）}

\noindent\begin{tabular}{@{}p{0.58\textwidth}@{\hspace{2em}}p{0.36\textwidth}@{}}
a.~纵声子（longitudinal phonons） & 固体；量子玻色液体；某些量子费米超流体\\
b.~横声子（transverse phonons） & 固体\\
c.~零声（zero sound） & 某些量子费米正常流体\\
d.~自旋密度波（spin density waves） & 某些量子费米正常流体\\
e.~等离激元（plasmons） & 金属和许多绝缘体——大多数带电的量子流体\\
\end{tabular}

（译注：2a 右栏原书印作“quantum Base liquid”,“Base”应为“Bose”。）

\emph{注意：}经典流体中的普通声和第二声（second sound）都\emph{不是}上述意义上的元激发，而是纯粹的统计现象；它们的存在，依赖于通过不可逆过程建立起局域平衡。

\emph{3. 自旋涨落（spin fluctuations）}

\noindent\begin{tabular}{@{}p{0.58\textwidth}@{\hspace{2em}}p{0.36\textwidth}@{}}
a.~铁磁自旋波：无能隙，$\Delta m=1$ & 铁磁体\\
b.~反铁磁自旋波：有能隙，当 $k\to 0$ 时 $\Delta m\to 0$ & 反铁磁体\\
c.~一般的自旋转动；$k$ 无定义，但仍是元激发理论的一个重要例子 & “磁状态”——铁磁体、反铁磁体和顺磁体\\
\end{tabular}

\emph{4. 其他（miscellaneous）}

\noindent\begin{tabular}{@{}p{0.58\textwidth}@{\hspace{2em}}p{0.36\textwidth}@{}}
a.~激子，包括 Slater 的自旋激发波 & 绝缘体的一种二粒子激发，或许还有超导体\\
b.~旋子（rotons）（有能隙，$k$ 空间中的一个面） & 液态 $\mathrm{He}_{4}$\\
c.~“d-mons” & （在含有两种或多种质量相差悬殊的载流子的物质中，一种其存在性尚成问题的密度波）\\
d.~与周期结构中的各种缺陷相联系的激发——表面态、局域模、杂质态、束缚激子等等。这些是研究细节的人的事情——它们对真实固体的行为可能至关重要，但在多体问题（m.b.p.）中却没有多少基本的重要性。 & \\
\end{tabular}

% ---- A 节收尾（原书 p.104 / PDF p119 顶部，协调者补译） ----
正如你们所见，固态物理学的多数有趣现象都可以在这张清单的某处找到自己的位置。然而并非全部；实际上，甚至并非所有那些与它们共享如下特征的存在物都列于其中：它们是正常固体的规则结构的小尺度、基元性的扰动，而我们在许多方面把它们当作独立粒子来处理。空位、间隙原子与杂质这类存在物，似乎应当另立一个“基元性”（elementaryness）的范畴，它与元激发的观念并不十分相似，主要在于它们的行为在真实的意义上是经典的而非量子力学的。位错，以及它们的“量子”对应物——量子化涡旋与磁通线——则又是另一个范畴；这些是尺度大得多（尺寸 $\sim N^{1/3}$）也更为复杂的现象。

对于每一种元激发类型，上面提到的近似独立性使我们能够用一个共同的进路去计算和理解热学与输运性质。作为零级近似，我们把元激发当作一种无相互作用的 Bose 或 Fermi 粒子气体（视情形而定），于是统计力学与输运理论就变得和完美无相互作用气体的情形一样容易。在更精细的应用中，我们可以把元激发之间的相互作用按任意想要的阶次包括进来；最低阶将允许它们在彼此接近到短距离时相互散射。散射元激发的相互作用，与元激发本身一样，必须被视为实际计算的结果——这种计算必然带入系统的 $N$ 体特征——但在许多情形下我们可以唯象地或近似地处理它，并得到合理的结果。最后，在许多重要情形下，我们甚至可以从元激发的性质出发反推关于基态本身性质的结论，声子零点运动的 Debye-Waller 理论就是一例。
